% Seminar 2 - Pagina 9. Fragment pentru % Seminar 2 - Pagina 9. Fragment pentru % Seminar 2 - Pagina 9. Fragment pentru % Seminar 2 - Pagina 9. Fragment pentru \input{seminars/pages/seminar2-pagina9-ro}.
% Fara preambul si fara intrerupere fortata de pagina.
% NOTE EDITORIALE:
% - Prima parte continua problema 11 de la pagina 8. In dezvoltarea
%   termenului final manuscrisul pare sa scrie (-1)^{n-2} C_n^{n-1},
%   desi suma generala care urmeaza impune (-1)^{n-1}. Am pastrat
%   exponentul vizibil si nu am corectat expresia tacit.
% - Problema despre deranjamente este numerotata 12. Urmeaza o problema
%   de corespondente Hom, al carei numar pare 13, iar demonstratia de
%   deranjamente din partea de jos are din nou numarul 12 in manuscris.
% - Exemplul de permutari pentru n=3 contine o a treia linie corectata;
%   sunt redate cele doua permutari lizibile nebarate.
% - La punctul a), o expresie initiala Hom(X,Y x Z) este taiata in
%   manuscris si apoi inlocuita cu Hom(X x Y,Z). Este redata a doua.
% - Ultima formula pentru deranjamente are limita superioara n-1, asa
%   cum pare scris. Formula completa standard ar mai include termenul
%   k=n. Am pastrat limita din manuscris, semnaland aici neconcordanta.
% - Desenele din dreapta sunt scheme pentru proiectiile produsului;
%   diagrama TikZ reprezinta relatiile lizibile, fara a completa alte probe.

\subsection*{Seminar 2 - Pagina 9}

\noindent\underline{Continuare - problema 11.} Dacă $n<m$, atunci
\[
\begin{aligned}
 s(m,n)={}&n^m-\binom{n}{1}(n-1)^m
          +\binom{n}{2}(n-2)^m-\binom{n}{3}(n-3)^m+\cdots\\
         &+(-1)^{n-2}\binom{n}{n-1}
          =\boxed{\displaystyle\sum_{k=0}^{n-1}(-1)^k\binom{n}{k}(n-k)^m}.
\end{aligned}
\]

\noindent\underline{Consecinţă:} pentru $m=n$,
\[
 \boxed{\displaystyle n!=\sum_{k=0}^{n}(-1)^k\binom{n}{k}(n-k)^n}.
\]

\bigskip
\noindent\fbox{12}\quad $\sigma\in S_n$ se numeşte \underline{deranjament}
 dacă $\sigma(i)\neq i$ pentru orice $i=\overline{1,n}$.

\noindent $\delta(n)$ - numărul deranjamentelor. Calculaţi $\delta(n)$.

\noindent\underline{Exemplu} ($n=3$):
\[
 \begin{pmatrix}1&2&3\\3&1&2\end{pmatrix},\qquad
 \begin{pmatrix}1&2&3\\2&3&1\end{pmatrix}.
\]

\medskip
\noindent\fbox{TEMĂ}

\medskip
\noindent\fbox{13}\quad Fie $X,Y,Z$ trei mulţimi. Atunci:
\[
 \text{a)}\quad
 \operatorname{Hom}(X\times Y,Z)
 \overset{\alpha}{\cong}
 \operatorname{Hom}\bigl(X,\operatorname{Hom}(Y,Z)\bigr),
\]
\[
 \text{b)}\quad
 \operatorname{Hom}(X,Y\times Z)
 \overset{\beta}{\cong}
 \operatorname{Hom}(X,Y)\times\operatorname{Hom}(X,Z).
\]
\noindent La b), aplicaţia este dată de
\[
 \beta(f)=\bigl(P_Y\circ f,\,P_Z\circ f\bigr).
\]
\[
\begin{tikzcd}[column sep=large,row sep=large]
 & X \arrow[dl,swap,"P_Y\circ f"] \arrow[d,"f"]
     \arrow[dr,"P_Z\circ f"] & \\
 Y & Y\times Z \arrow[l,"P_Y"] \arrow[r,swap,"P_Z"] & Z
\end{tikzcd}
\]
\noindent $P_Y(y,z)=y$, $P_Z(y,z)=z$.

\bigskip
\noindent\fbox{12}\quad \underline{Calculul deranjamentelor (continuare).}

\noindent $\delta(n)=n!$ minus numărul permutărilor pentru care
 există $i$ a.î. $\sigma(i)=i$.
\[
 A_i=\{\sigma\in S_n\mid\sigma(i)=i\}.
\]
\noindent Avem
\[
 |A_i|=(n-1)!,\qquad
 |A_i\cap A_j|=(n-2)!,\qquad
 |A_i\cap A_j\cap A_k|=(n-3)!.
\]
\[
 \delta_n=n!-\bigl|A_1\cup A_2\cup\cdots\cup A_n\bigr|.
\]
\noindent Prin principiul includerii şi excluderii,
\[
\begin{aligned}
 \bigl|A_1\cup\cdots\cup A_n\bigr|
  ={}&n(n-1)!-\binom{n}{2}(n-2)!+\binom{n}{3}(n-3)!+\cdots\\
     &+(-1)^{n-2}\binom{n}{n-1}\cdot1\\
  ={}&\sum_{k=1}^{n-1}\binom{n}{k}(-1)^{k-1}(n-k)!.
\end{aligned}
\]
\noindent Rezultă, în notaţia manuscrisului,
\[
\begin{aligned}
 \delta_n
  &=n!-\sum_{k=1}^{n-1}(-1)^{k-1}\binom{n}{k}(n-k)!\\
  &=\boxed{\displaystyle\sum_{k=0}^{n-1}(-1)^k\binom{n}{k}(n-k)!}.
\end{aligned}
\]
.
% Fara preambul si fara intrerupere fortata de pagina.
% NOTE EDITORIALE:
% - Prima parte continua problema 11 de la pagina 8. In dezvoltarea
%   termenului final manuscrisul pare sa scrie (-1)^{n-2} C_n^{n-1},
%   desi suma generala care urmeaza impune (-1)^{n-1}. Am pastrat
%   exponentul vizibil si nu am corectat expresia tacit.
% - Problema despre deranjamente este numerotata 12. Urmeaza o problema
%   de corespondente Hom, al carei numar pare 13, iar demonstratia de
%   deranjamente din partea de jos are din nou numarul 12 in manuscris.
% - Exemplul de permutari pentru n=3 contine o a treia linie corectata;
%   sunt redate cele doua permutari lizibile nebarate.
% - La punctul a), o expresie initiala Hom(X,Y x Z) este taiata in
%   manuscris si apoi inlocuita cu Hom(X x Y,Z). Este redata a doua.
% - Ultima formula pentru deranjamente are limita superioara n-1, asa
%   cum pare scris. Formula completa standard ar mai include termenul
%   k=n. Am pastrat limita din manuscris, semnaland aici neconcordanta.
% - Desenele din dreapta sunt scheme pentru proiectiile produsului;
%   diagrama TikZ reprezinta relatiile lizibile, fara a completa alte probe.

\subsection*{Seminar 2 - Pagina 9}

\noindent\underline{Continuare - problema 11.} Dacă $n<m$, atunci
\[
\begin{aligned}
 s(m,n)={}&n^m-\binom{n}{1}(n-1)^m
          +\binom{n}{2}(n-2)^m-\binom{n}{3}(n-3)^m+\cdots\\
         &+(-1)^{n-2}\binom{n}{n-1}
          =\boxed{\displaystyle\sum_{k=0}^{n-1}(-1)^k\binom{n}{k}(n-k)^m}.
\end{aligned}
\]

\noindent\underline{Consecinţă:} pentru $m=n$,
\[
 \boxed{\displaystyle n!=\sum_{k=0}^{n}(-1)^k\binom{n}{k}(n-k)^n}.
\]

\bigskip
\noindent\fbox{12}\quad $\sigma\in S_n$ se numeşte \underline{deranjament}
 dacă $\sigma(i)\neq i$ pentru orice $i=\overline{1,n}$.

\noindent $\delta(n)$ - numărul deranjamentelor. Calculaţi $\delta(n)$.

\noindent\underline{Exemplu} ($n=3$):
\[
 \begin{pmatrix}1&2&3\\3&1&2\end{pmatrix},\qquad
 \begin{pmatrix}1&2&3\\2&3&1\end{pmatrix}.
\]

\medskip
\noindent\fbox{TEMĂ}

\medskip
\noindent\fbox{13}\quad Fie $X,Y,Z$ trei mulţimi. Atunci:
\[
 \text{a)}\quad
 \operatorname{Hom}(X\times Y,Z)
 \overset{\alpha}{\cong}
 \operatorname{Hom}\bigl(X,\operatorname{Hom}(Y,Z)\bigr),
\]
\[
 \text{b)}\quad
 \operatorname{Hom}(X,Y\times Z)
 \overset{\beta}{\cong}
 \operatorname{Hom}(X,Y)\times\operatorname{Hom}(X,Z).
\]
\noindent La b), aplicaţia este dată de
\[
 \beta(f)=\bigl(P_Y\circ f,\,P_Z\circ f\bigr).
\]
\[
\begin{tikzcd}[column sep=large,row sep=large]
 & X \arrow[dl,swap,"P_Y\circ f"] \arrow[d,"f"]
     \arrow[dr,"P_Z\circ f"] & \\
 Y & Y\times Z \arrow[l,"P_Y"] \arrow[r,swap,"P_Z"] & Z
\end{tikzcd}
\]
\noindent $P_Y(y,z)=y$, $P_Z(y,z)=z$.

\bigskip
\noindent\fbox{12}\quad \underline{Calculul deranjamentelor (continuare).}

\noindent $\delta(n)=n!$ minus numărul permutărilor pentru care
 există $i$ a.î. $\sigma(i)=i$.
\[
 A_i=\{\sigma\in S_n\mid\sigma(i)=i\}.
\]
\noindent Avem
\[
 |A_i|=(n-1)!,\qquad
 |A_i\cap A_j|=(n-2)!,\qquad
 |A_i\cap A_j\cap A_k|=(n-3)!.
\]
\[
 \delta_n=n!-\bigl|A_1\cup A_2\cup\cdots\cup A_n\bigr|.
\]
\noindent Prin principiul includerii şi excluderii,
\[
\begin{aligned}
 \bigl|A_1\cup\cdots\cup A_n\bigr|
  ={}&n(n-1)!-\binom{n}{2}(n-2)!+\binom{n}{3}(n-3)!+\cdots\\
     &+(-1)^{n-2}\binom{n}{n-1}\cdot1\\
  ={}&\sum_{k=1}^{n-1}\binom{n}{k}(-1)^{k-1}(n-k)!.
\end{aligned}
\]
\noindent Rezultă, în notaţia manuscrisului,
\[
\begin{aligned}
 \delta_n
  &=n!-\sum_{k=1}^{n-1}(-1)^{k-1}\binom{n}{k}(n-k)!\\
  &=\boxed{\displaystyle\sum_{k=0}^{n-1}(-1)^k\binom{n}{k}(n-k)!}.
\end{aligned}
\]
.
% Fara preambul si fara intrerupere fortata de pagina.
% NOTE EDITORIALE:
% - Prima parte continua problema 11 de la pagina 8. In dezvoltarea
%   termenului final manuscrisul pare sa scrie (-1)^{n-2} C_n^{n-1},
%   desi suma generala care urmeaza impune (-1)^{n-1}. Am pastrat
%   exponentul vizibil si nu am corectat expresia tacit.
% - Problema despre deranjamente este numerotata 12. Urmeaza o problema
%   de corespondente Hom, al carei numar pare 13, iar demonstratia de
%   deranjamente din partea de jos are din nou numarul 12 in manuscris.
% - Exemplul de permutari pentru n=3 contine o a treia linie corectata;
%   sunt redate cele doua permutari lizibile nebarate.
% - La punctul a), o expresie initiala Hom(X,Y x Z) este taiata in
%   manuscris si apoi inlocuita cu Hom(X x Y,Z). Este redata a doua.
% - Ultima formula pentru deranjamente are limita superioara n-1, asa
%   cum pare scris. Formula completa standard ar mai include termenul
%   k=n. Am pastrat limita din manuscris, semnaland aici neconcordanta.
% - Desenele din dreapta sunt scheme pentru proiectiile produsului;
%   diagrama TikZ reprezinta relatiile lizibile, fara a completa alte probe.

\subsection*{Seminar 2 - Pagina 9}

\noindent\underline{Continuare - problema 11.} Dacă $n<m$, atunci
\[
\begin{aligned}
 s(m,n)={}&n^m-\binom{n}{1}(n-1)^m
          +\binom{n}{2}(n-2)^m-\binom{n}{3}(n-3)^m+\cdots\\
         &+(-1)^{n-2}\binom{n}{n-1}
          =\boxed{\displaystyle\sum_{k=0}^{n-1}(-1)^k\binom{n}{k}(n-k)^m}.
\end{aligned}
\]

\noindent\underline{Consecinţă:} pentru $m=n$,
\[
 \boxed{\displaystyle n!=\sum_{k=0}^{n}(-1)^k\binom{n}{k}(n-k)^n}.
\]

\bigskip
\noindent\fbox{12}\quad $\sigma\in S_n$ se numeşte \underline{deranjament}
 dacă $\sigma(i)\neq i$ pentru orice $i=\overline{1,n}$.

\noindent $\delta(n)$ - numărul deranjamentelor. Calculaţi $\delta(n)$.

\noindent\underline{Exemplu} ($n=3$):
\[
 \begin{pmatrix}1&2&3\\3&1&2\end{pmatrix},\qquad
 \begin{pmatrix}1&2&3\\2&3&1\end{pmatrix}.
\]

\medskip
\noindent\fbox{TEMĂ}

\medskip
\noindent\fbox{13}\quad Fie $X,Y,Z$ trei mulţimi. Atunci:
\[
 \text{a)}\quad
 \operatorname{Hom}(X\times Y,Z)
 \overset{\alpha}{\cong}
 \operatorname{Hom}\bigl(X,\operatorname{Hom}(Y,Z)\bigr),
\]
\[
 \text{b)}\quad
 \operatorname{Hom}(X,Y\times Z)
 \overset{\beta}{\cong}
 \operatorname{Hom}(X,Y)\times\operatorname{Hom}(X,Z).
\]
\noindent La b), aplicaţia este dată de
\[
 \beta(f)=\bigl(P_Y\circ f,\,P_Z\circ f\bigr).
\]
\[
\begin{tikzcd}[column sep=large,row sep=large]
 & X \arrow[dl,swap,"P_Y\circ f"] \arrow[d,"f"]
     \arrow[dr,"P_Z\circ f"] & \\
 Y & Y\times Z \arrow[l,"P_Y"] \arrow[r,swap,"P_Z"] & Z
\end{tikzcd}
\]
\noindent $P_Y(y,z)=y$, $P_Z(y,z)=z$.

\bigskip
\noindent\fbox{12}\quad \underline{Calculul deranjamentelor (continuare).}

\noindent $\delta(n)=n!$ minus numărul permutărilor pentru care
 există $i$ a.î. $\sigma(i)=i$.
\[
 A_i=\{\sigma\in S_n\mid\sigma(i)=i\}.
\]
\noindent Avem
\[
 |A_i|=(n-1)!,\qquad
 |A_i\cap A_j|=(n-2)!,\qquad
 |A_i\cap A_j\cap A_k|=(n-3)!.
\]
\[
 \delta_n=n!-\bigl|A_1\cup A_2\cup\cdots\cup A_n\bigr|.
\]
\noindent Prin principiul includerii şi excluderii,
\[
\begin{aligned}
 \bigl|A_1\cup\cdots\cup A_n\bigr|
  ={}&n(n-1)!-\binom{n}{2}(n-2)!+\binom{n}{3}(n-3)!+\cdots\\
     &+(-1)^{n-2}\binom{n}{n-1}\cdot1\\
  ={}&\sum_{k=1}^{n-1}\binom{n}{k}(-1)^{k-1}(n-k)!.
\end{aligned}
\]
\noindent Rezultă, în notaţia manuscrisului,
\[
\begin{aligned}
 \delta_n
  &=n!-\sum_{k=1}^{n-1}(-1)^{k-1}\binom{n}{k}(n-k)!\\
  &=\boxed{\displaystyle\sum_{k=0}^{n-1}(-1)^k\binom{n}{k}(n-k)!}.
\end{aligned}
\]
.
% Fara preambul si fara intrerupere fortata de pagina.
% NOTE EDITORIALE:
% - Prima parte continua problema 11 de la pagina 8. In dezvoltarea
%   termenului final manuscrisul pare sa scrie (-1)^{n-2} C_n^{n-1},
%   desi suma generala care urmeaza impune (-1)^{n-1}. Am pastrat
%   exponentul vizibil si nu am corectat expresia tacit.
% - Problema despre deranjamente este numerotata 12. Urmeaza o problema
%   de corespondente Hom, al carei numar pare 13, iar demonstratia de
%   deranjamente din partea de jos are din nou numarul 12 in manuscris.
% - Exemplul de permutari pentru n=3 contine o a treia linie corectata;
%   sunt redate cele doua permutari lizibile nebarate.
% - La punctul a), o expresie initiala Hom(X,Y x Z) este taiata in
%   manuscris si apoi inlocuita cu Hom(X x Y,Z). Este redata a doua.
% - Ultima formula pentru deranjamente are limita superioara n-1, asa
%   cum pare scris. Formula completa standard ar mai include termenul
%   k=n. Am pastrat limita din manuscris, semnaland aici neconcordanta.
% - Desenele din dreapta sunt scheme pentru proiectiile produsului;
%   diagrama TikZ reprezinta relatiile lizibile, fara a completa alte probe.

\subsection*{Seminar 2 - Pagina 9}

\noindent\underline{Continuare - problema 11.} Dacă $n<m$, atunci
\[
\begin{aligned}
 s(m,n)={}&n^m-\binom{n}{1}(n-1)^m
          +\binom{n}{2}(n-2)^m-\binom{n}{3}(n-3)^m+\cdots\\
         &+(-1)^{n-2}\binom{n}{n-1}
          =\boxed{\displaystyle\sum_{k=0}^{n-1}(-1)^k\binom{n}{k}(n-k)^m}.
\end{aligned}
\]

\noindent\underline{Consecinţă:} pentru $m=n$,
\[
 \boxed{\displaystyle n!=\sum_{k=0}^{n}(-1)^k\binom{n}{k}(n-k)^n}.
\]

\bigskip
\noindent\fbox{12}\quad $\sigma\in S_n$ se numeşte \underline{deranjament}
 dacă $\sigma(i)\neq i$ pentru orice $i=\overline{1,n}$.

\noindent $\delta(n)$ - numărul deranjamentelor. Calculaţi $\delta(n)$.

\noindent\underline{Exemplu} ($n=3$):
\[
 \begin{pmatrix}1&2&3\\3&1&2\end{pmatrix},\qquad
 \begin{pmatrix}1&2&3\\2&3&1\end{pmatrix}.
\]

\medskip
\noindent\fbox{TEMĂ}

\medskip
\noindent\fbox{13}\quad Fie $X,Y,Z$ trei mulţimi. Atunci:
\[
 \text{a)}\quad
 \operatorname{Hom}(X\times Y,Z)
 \overset{\alpha}{\cong}
 \operatorname{Hom}\bigl(X,\operatorname{Hom}(Y,Z)\bigr),
\]
\[
 \text{b)}\quad
 \operatorname{Hom}(X,Y\times Z)
 \overset{\beta}{\cong}
 \operatorname{Hom}(X,Y)\times\operatorname{Hom}(X,Z).
\]
\noindent La b), aplicaţia este dată de
\[
 \beta(f)=\bigl(P_Y\circ f,\,P_Z\circ f\bigr).
\]
\[
\begin{tikzcd}[column sep=large,row sep=large]
 & X \arrow[dl,swap,"P_Y\circ f"] \arrow[d,"f"]
     \arrow[dr,"P_Z\circ f"] & \\
 Y & Y\times Z \arrow[l,"P_Y"] \arrow[r,swap,"P_Z"] & Z
\end{tikzcd}
\]
\noindent $P_Y(y,z)=y$, $P_Z(y,z)=z$.

\bigskip
\noindent\fbox{12}\quad \underline{Calculul deranjamentelor (continuare).}

\noindent $\delta(n)=n!$ minus numărul permutărilor pentru care
 există $i$ a.î. $\sigma(i)=i$.
\[
 A_i=\{\sigma\in S_n\mid\sigma(i)=i\}.
\]
\noindent Avem
\[
 |A_i|=(n-1)!,\qquad
 |A_i\cap A_j|=(n-2)!,\qquad
 |A_i\cap A_j\cap A_k|=(n-3)!.
\]
\[
 \delta_n=n!-\bigl|A_1\cup A_2\cup\cdots\cup A_n\bigr|.
\]
\noindent Prin principiul includerii şi excluderii,
\[
\begin{aligned}
 \bigl|A_1\cup\cdots\cup A_n\bigr|
  ={}&n(n-1)!-\binom{n}{2}(n-2)!+\binom{n}{3}(n-3)!+\cdots\\
     &+(-1)^{n-2}\binom{n}{n-1}\cdot1\\
  ={}&\sum_{k=1}^{n-1}\binom{n}{k}(-1)^{k-1}(n-k)!.
\end{aligned}
\]
\noindent Rezultă, în notaţia manuscrisului,
\[
\begin{aligned}
 \delta_n
  &=n!-\sum_{k=1}^{n-1}(-1)^{k-1}\binom{n}{k}(n-k)!\\
  &=\boxed{\displaystyle\sum_{k=0}^{n-1}(-1)^k\binom{n}{k}(n-k)!}.
\end{aligned}
\]
